% =====================================================================
% conclusao.tex — Conclusão e Trabalhos Futuros
% =====================================================================

\chapter{Conclusão}
\label{cap:conclusao}

\section{Síntese das Contribuições}

Este Projeto Final de Curso apresentou o desenvolvimento integral de
um sistema de Caixão de Areia com Realidade Aumentada para apoio ao
treinamento militar em operações no terreno. O sistema combina,
de maneira engenhosa, três elementos centrais: (i) o rigor matemático
de uma calibração geométrica completa baseada em SVD, Gram-Schmidt e
modelo \textit{pinhole} de Tsai; (ii) a regra de negócio quantitativa
de comparação célula a célula contra um Modelo Digital de Elevação
fornecido em formato GeoTIFF, com tolerância vertical explícita; e
(iii) uma arquitetura resiliente, dotada de cadeia de
\textit{fallback} em quatro níveis, que assegura execução
\textit{plug and play} mesmo na ausência total de hardware ou de
dados cartográficos.

O motor matemático foi desenvolvido sob a metodologia de
\textit{Test-Driven Development}, resultando em uma suíte com 88
testes unitários automatizados que verificam — com valores numéricos
explícitos e auditáveis — todos os contratos algébricos do sistema,
da calibração RANSAC + SVD ao teste de aceitação de ponta a ponta.
A arquitetura monolítica modularizada, com seus quatro padrões de
projeto (Strategy + Fallback, Adapter, State Machine, Observer),
garante simultaneamente a latência mínima necessária para operação
em tempo real e a separação clara de responsabilidades indispensável
à manutenção.

A renderização por malha discretizada — em substituição à projeção
ingênua de pontos isolados — produz \textit{feedback} visual
contínuo, geometricamente preciso, computacionalmente eficiente e
estatisticamente robusto ao ruído do sensor. O emulador interativo
controlado por \emph{mouse} permite, em qualquer máquina sem hardware
especializado, demonstrar todo o \textit{pipeline} de forma tangível,
funcionalidade especialmente valiosa para defesa do projeto e para
ensino em sala de aula.

Durante a fase de validação em campo com o sensor Kinect~v2 real, foi
identificado e corrigido um defeito de calibração que reduzia a área
renderizada a uma pequena fração da janela de projeção. A correção,
formalizada na composição com a matriz $T_{\text{shift}}$ e na
introdução do filtro de extensão da mesa, foi documentada com
detalhes no Capítulo~\ref{cap:resultados} e permanece consolidada nas
versões subsequentes do sistema, que evoluíram a calibração para o
esquema RANSAC~+~SVD sobre uma superfície de referência (tampa ou
base vazia) com cota zero fixada na base do caixão.

\section{Atendimento aos Objetivos}

Todos os oito objetivos específicos enunciados no
Capítulo~\ref{cap:intro} foram atendidos:

\begin{enumerate}
    \item Arquitetura de três camadas implementada e documentada;
    \item Cadeia de \textit{fallback} do \texttt{KinectSensor} operacional;
    \item Motor matemático com SVD + Gram-Schmidt + Tsai funcional;
    \item \texttt{AdaptadorMDE} implementado, com \emph{fallback} gaussiano
          validado por teste automatizado; a leitura de GeoTIFF real está
          implementada e integrada ao restante do \textit{pipeline}, mas
          — como registrado com transparência na
          Seção~\ref{sec:tdd} — ainda não possui um caso de teste
          automatizado dedicado nem foi demonstrada com o arquivo
          \texttt{25S51\_ZN.tif} versionado no repositório, permanecendo
          como item de verificação futura (Seção~\ref{sec:trabalhos-futuros});
    \item Discretização em malha com renderização por \texttt{cv2.fillPoly};
    \item Emulador interativo via \texttt{cv2.setMouseCallback};
    \item Suíte de 88 testes unitários, todos aprovados;
    \item Documentação tripla (acadêmica, prática e relatório ABNT).
\end{enumerate}

\section{Trabalhos Futuros}
\label{sec:trabalhos-futuros}

A continuidade da linha de pesquisa pode explorar, dentre outras, as
seguintes frentes:

\begin{description}
    \item[Teste automatizado de leitura de GeoTIFF real.] A suíte atual
          valida a leitura de MDE apenas contra as superfícies
          sintéticas Cubo Central e Morro Gaussiano (Seção~\ref{sec:tdd}).
          Escrever um teste dedicado que carregue o GeoTIFF real já
          versionado no repositório (\texttt{25S51\_ZN.tif}) e valide a
          interpolação bilinear contra valores conhecidos da grade
          fecharia essa lacuna de cobertura.
    \item[Validação do erro quantitativo do MDE com ruído real de
          sensor.] O erro medido na Seção~\ref{subsec:erro-quantitativo-mde}
          (RMSE de 11,4\%) usa uma nuvem sintética livre de ruído — o
          único erro presente é a quantização geométrica da grade. O
          Kinect real introduz ruído próprio de profundidade,
          tipicamente mais acentuado nas bordas de objetos, que essa
          medição não captura; repetir a mesma análise com nuvens
          capturadas por hardware real, e para o mapa Morro Gaussiano
          além do Cubo Central, forneceria uma cota superior de erro
          mais realista.
    \item[Reconstrução do caixão de teste e nova demonstração física
          integral.] Corrigir o recorte da estrutura de madeira que
          hoje tampa parcialmente o campo de visão do sensor
          (Seção~\ref{subsec:limitacao-fisica}) e repetir o roteiro de
          aceitação completo com Kinect real, do início ao fim, sem
          interrupções de conexão.
    \item[Benchmark de desempenho em modo real.] Medir e reportar a
          taxa de quadros com o sensor Kinect conectado, de forma
          independente da medição já feita em modo simulação
          (Seção~\ref{sec:funcional}).
    \item[Calibração automática do projetor.] A versão atual emprega
          parâmetros de projeção configurados analiticamente. Uma
          calibração genuína do projetor por padrões \textit{Gray code}
          ou tabuleiro de xadrez (já parcialmente suportada pelas
          funções \texttt{encontrar\_\allowbreak cantos\_\allowbreak
          tabuleiro} e \texttt{calibrar\_\allowbreak projetor})
          eliminaria distorções residuais em montagens não-ideais.
    \item[Suporte a múltiplas camadas de visualização.] Sobreposição
          de curvas de nível, escalas batimétricas, isótermas e
          itinerários táticos — todos coerentes com o mesmo MDE de
          referência.
    \item[Integração com SARAS-FFT (Águas Pluviais).] Simulação
          hidrológica de escoamento sobre o relevo modelado, com
          renderização em tempo real do fluxo, em complemento à
          modelagem topográfica.
    \item[Substituição do Kinect por sensor moderno.] O Microsoft
          Kinect foi descontinuado; alternativas como o
          Intel RealSense, o Stereolabs ZED e o
          Microsoft Azure Kinect DK oferecem maior
          resolução e confiabilidade. A arquitetura
          \textit{strategy + fallback} já está preparada para
          acomodar essa migração com mínimo impacto.
    \item[Validação em escola militar.] Estudo experimental
          comparando o tempo médio de modelagem do terreno por
          turmas com e sem o sistema, e a fidelidade cartográfica
          do produto final medida por reconstrução fotogramétrica
          independente.
    \item[Exportação para realidade virtual.] A nuvem de pontos
          calibrada pode alimentar um \textit{viewer} VR (Oculus,
          HTC Vive) para imersão remota ou de instrução à distância.
\end{description}

\section{Considerações Finais}

A combinação de matemática rigorosa, engenharia de software
disciplinada e visão clara da aplicação militar produziu um sistema
que se diferencia das implementações de caixão de areia com realidade
aumentada de finalidade puramente didática ou recreativa levantadas na
revisão de trabalhos relacionados (Capítulo~\ref{cap:relacionados}),
por incorporar comparação quantitativa contra um MDE de referência,
tolerância vertical explícita e uma arquitetura de \textit{fallback}
resiliente. A entrega não se limita a um \textit{software} funcional: contempla
testes automatizados, documentação acadêmica completa, guia de uso
prático e este próprio relatório de conformidade ABNT, formando um
conjunto reproduzível, auditável e diretamente aproveitável pelas
escolas de formação militar.

Acreditamos que o trabalho contribui, ainda que modestamente, para
a tradição de excelência técnica do Instituto Militar de
Engenharia, e esperamos que sirva de ponto de partida para futuras
melhorias e adaptações por gerações seguintes de cadetes e
engenheiros militares.
